\documentclass[12pt,a4paper]{amsart}

\usepackage{amsmath,amssymb,amsthm}
\usepackage{mathtools}
\usepackage{hyperref}
\usepackage{geometry}
\geometry{margin=2.5cm}

% -------------------------------------------------------
% Theorem environments
% -------------------------------------------------------
\newtheorem{theorem}{Theorem}[section]
\newtheorem{lemma}[theorem]{Lemma}
\newtheorem{corollary}[theorem]{Corollary}
\newtheorem{proposition}[theorem]{Proposition}
\theoremstyle{definition}
\newtheorem{definition}[theorem]{Definition}
\newtheorem{remark}[theorem]{Remark}

% -------------------------------------------------------
% Custom commands
% -------------------------------------------------------
\newcommand{\N}{\mathbb{N}}
\newcommand{\Z}{\mathbb{Z}}
\newcommand{\Q}{\mathbb{Q}}
\newcommand{\Fp}{\mathbb{F}_p}
\DeclareMathOperator{\ord}{ord}
\newcommand{\leg}[2]{\left(\frac{#1}{#2}\right)}

% -------------------------------------------------------
\begin{document}
% -------------------------------------------------------

\title{Applications of Wilson's Theorem}

\author{Michael Fuhrmann}
\address{Specialist Mathematics Project, Build 66}
\email{specialist-maths@localhost}
\date{\today}

\begin{abstract}
We demonstrate four major applications of Wilson's theorem:
(i) a complete primality characterisation via $(n-1)! \equiv -1 \pmod n$,
with an analysis of its computational impracticality;
(ii) a proof that $-1$ is a quadratic residue modulo an odd prime $p$
if and only if $p \equiv 1 \pmod 4$, derived directly from Wilson via
the factorisation of $(p-1)!$;
(iii) an elementary proof of Wolstenholme's harmonic lemma:
$1 + 1/2 + \cdots + 1/(p-1) \equiv 0 \pmod p$ for all primes $p \geq 3$;
and (iv) an alternative proof of Fermat's little theorem using Wilson's theorem,
with a sketch of how this approach extends to Euler's theorem.
All proofs are elementary and self-contained.
\end{abstract}

\maketitle
\tableofcontents

% ================================================================
\section{Introduction}
% ================================================================

Wilson's theorem — $(p-1)! \equiv -1 \pmod p$ if and only if $p$ is prime —
is more than an isolated curiosity.  It is a tool for deriving deeper facts
about primes, quadratic residues, and multiplicative arithmetic.

In this paper we collect four applications.  The first is the theorem itself
used as a primality test, with an honest assessment of its limitations.
The second derives a complete criterion for $-1$ to be a quadratic residue
modulo $p$ — a result usually proved via character theory or Gauss sums — here
obtained with nothing more than high school algebra.  The third proves a
harmonic congruence due to Wolstenholme.  The fourth gives an alternative
proof of Fermat's little theorem via Wilson, demonstrating the interconnection
of these classical results.

% ================================================================
\section{Wilson's Theorem as a Primality Test}
% ================================================================

\begin{theorem}[Primality characterisation]\label{thm:primality}
  An integer $n \geq 2$ is prime if and only if
  \[
    (n-1)! \;\equiv\; -1 \pmod n.
  \]
\end{theorem}

\begin{proof}
  \textbf{($\Rightarrow$)} Suppose $n = p$ is prime.  The multiplicative group
  $(\Z/p\Z)^*$ has order $p-1$, and its elements are $1, 2, \ldots, p-1$.
  Every element $a$ has a unique inverse $a^{-1}$ in $\{1, \ldots, p-1\}$.
  The only elements equal to their own inverse satisfy $a^2 \equiv 1 \pmod p$,
  i.e.\ $p \mid (a-1)(a+1)$; since $p$ is prime, $a \equiv \pm 1 \pmod p$,
  giving $a \in \{1, p-1\}$.

  All other elements $2, \ldots, p-2$ pair up: $a \leftrightarrow a^{-1}$
  with $a \cdot a^{-1} \equiv 1 \pmod p$.  Hence
  \[
    (p-1)! = 1 \cdot (p-1) \cdot \prod_{\substack{a=2\\a \neq a^{-1}}}^{p-2} a
    \equiv 1 \cdot (p-1) \cdot 1 \equiv p-1 \equiv -1 \pmod p.
  \]

  \textbf{($\Leftarrow$)} Suppose $n$ is composite, $n = ab$ with $1 < a \leq b < n$.

  \emph{Case 1: $a < b$.}  Then $a$ and $b$ are distinct integers in
  $\{1, \ldots, n-1\}$, so $ab = n \mid (n-1)!$ and $(n-1)! \equiv 0 \not\equiv -1$.

  \emph{Case 2: $a = b$, so $n = a^2$.}
  If $a \geq 3$: both $a$ and $2a$ lie in $\{1, \ldots, n-1\}$ (since
  $2a < a^2$ for $a \geq 3$), so $2a^2 = 2n \mid (n-1)!$ and $(n-1)! \equiv 0$.
  If $a = 2$: $n = 4$ and $(4-1)! = 6 \equiv 2 \pmod 4 \neq -1$.
\end{proof}

\begin{remark}[Computational impracticality]
  Although Theorem~\ref{thm:primality} provides a perfect primality test in
  principle, it is computationally useless for large $n$.  Computing $(n-1)!
  \bmod n$ naively requires $n-2$ multiplications, each modulo $n$.
  With $n \approx 2^{1000}$ (a $300$-digit prime), this needs $\sim 2^{1000}$
  steps — far beyond any feasible computation.

  By contrast, Miller--Rabin probabilistic primality testing runs in
  $O(k \log^2 n)$ time with $k$ rounds, and the AKS algorithm \cite{Agrawal2004}
  deterministically decides primality in polynomial time $O(\log^{6} n)$.
  Wilson's test has theoretical importance but no practical application to
  primality testing.
\end{remark}

% ================================================================
\section{Quadratic Residues: When Is $-1$ a Square Modulo $p$?}
% ================================================================

\begin{definition}[Quadratic residue]\label{def:quadratic_residue}
  Let $p$ be an odd prime.  An integer $a$ with $\gcd(a,p) = 1$ is a
  \emph{quadratic residue} modulo $p$ if $a \equiv x^2 \pmod p$ for some
  integer $x$.  The Legendre symbol is defined by
  \[
    \leg{a}{p} =
    \begin{cases}
      1  & \text{if } a \text{ is a quadratic residue mod } p, \\
      -1 & \text{if } a \text{ is a non-residue mod } p.
    \end{cases}
  \]
\end{definition}

\begin{theorem}[First supplement to quadratic reciprocity]\label{thm:minus1_qr}
  Let $p$ be an odd prime.  Then $-1$ is a quadratic residue modulo $p$
  if and only if $p \equiv 1 \pmod 4$.  Equivalently:
  \[
    \leg{-1}{p} \;=\; (-1)^{(p-1)/2}.
  \]
\end{theorem}

\begin{proof}
  We compute $(p-1)!$ in two different ways.

  \textbf{Way 1 (Wilson):}
  By Wilson's theorem, $(p-1)! \equiv -1 \pmod p$.

  \textbf{Way 2 (regrouping):}
  Write the product $(p-1)! = 1 \cdot 2 \cdots \frac{p-1}{2} \cdot \frac{p+1}{2}
  \cdots (p-1)$.  In the second half, replace $j$ by $p - j$ for
  $j \in \{(p+1)/2, \ldots, p-1\}$, which maps bijectively to
  $\{1, \ldots, (p-1)/2\}$ via $j \mapsto p-j$:
  \[
    \frac{p+1}{2} \cdot \frac{p+3}{2} \cdots (p-1)
    = (p - \tfrac{p-1}{2})(p - \tfrac{p-3}{2}) \cdots (p-1).
  \]
  More cleanly: for $k = 1, 2, \ldots, (p-1)/2$ we have
  $p - k \equiv -k \pmod p$.  So
  \[
    (p-1)! = \prod_{k=1}^{(p-1)/2} k \cdot \prod_{k=1}^{(p-1)/2}(p-k)
    \equiv \prod_{k=1}^{(p-1)/2} k \cdot \prod_{k=1}^{(p-1)/2}(-k)
    = \Bigl(\prod_{k=1}^{(p-1)/2} k\Bigr)^2 \cdot (-1)^{(p-1)/2} \pmod p.
  \]
  Let $M = \bigl(\tfrac{p-1}{2}\bigr)!$.  Then
  \[
    (p-1)! \equiv M^2 \cdot (-1)^{(p-1)/2} \pmod p.
  \]

  \textbf{Combining:}
  From Ways 1 and 2:
  \[
    -1 \equiv M^2 \cdot (-1)^{(p-1)/2} \pmod p.
  \]
  Therefore
  \[
    (-1)^{(p-1)/2} \equiv -1 \cdot M^{-2} \pmod p
    \quad\Longleftrightarrow\quad
    (-1)^{(p-1)/2+1} \equiv M^{-2} \pmod p.
  \]
  Rearranging:
  \[
    (-1)^{(p-1)/2} \equiv \frac{-1}{M^2} \pmod p.
  \]
  But we want to determine whether $-1$ is a QR.  From the equation
  $-1 \equiv M^2 (-1)^{(p-1)/2}$:
  \begin{itemize}
    \item If $(p-1)/2$ is even (i.e.\ $p \equiv 1 \pmod 4$):
      $(-1)^{(p-1)/2} = 1$, so $-1 \equiv M^2 \pmod p$.
      Thus $-1$ is a quadratic residue modulo $p$. \checkmark
    \item If $(p-1)/2$ is odd (i.e.\ $p \equiv 3 \pmod 4$):
      $(-1)^{(p-1)/2} = -1$, so $-1 \equiv -M^2 \pmod p$,
      i.e.\ $M^2 \equiv 1 \pmod p$, meaning $M \equiv \pm 1 \pmod p$.
      If $-1$ were a QR, then $-1 \equiv x^2 \pmod p$ for some $x$.
      Multiplying: $(-1)^{(p-1)/2} \equiv x^{p-1} \equiv 1 \pmod p$ by
      Fermat's little theorem.  But $(-1)^{(p-1)/2} = -1 \neq 1$. Contradiction.
      Hence $-1$ is a non-residue. \checkmark
  \end{itemize}
  In both cases the formula $\leg{-1}{p} = (-1)^{(p-1)/2}$ holds.
\end{proof}

\begin{corollary}[Fermat's two-square theorem — prerequisite]\label{cor:twosquares}
  A prime $p$ can be written as $p = x^2 + y^2$ with integers $x, y$ only
  if $p = 2$ or $p \equiv 1 \pmod 4$.
\end{corollary}

\begin{proof}
  If $p = x^2 + y^2$, then $x^2 \equiv -y^2 \pmod p$.  If $p \nmid y$, then
  $(xy^{-1})^2 \equiv -1 \pmod p$, so $-1$ is a QR mod $p$, requiring
  $p \equiv 1 \pmod 4$ by Theorem~\ref{thm:minus1_qr}.
  The case $p = 2 = 1^2 + 1^2$ is direct.
\end{proof}

% ================================================================
\section{Wolstenholme's Harmonic Lemma}
% ================================================================

\begin{theorem}[Wolstenholme's harmonic congruence]\label{thm:harmonic}
  For every prime $p \geq 3$:
  \[
    1 + \frac{1}{2} + \frac{1}{3} + \cdots + \frac{1}{p-1}
    \;\equiv\; 0 \pmod p,
  \]
  where each fraction $1/j$ denotes the multiplicative inverse of $j$ modulo $p$.
\end{theorem}

\begin{proof}
  Denote the harmonic sum $H = \sum_{j=1}^{p-1} j^{-1} \pmod p$.

  \textbf{Step 1 (Symmetry).}
  For each $j \in \{1, \ldots, p-1\}$, its inverse $j^{-1}$ also lies in
  $\{1, \ldots, p-1\}$ (since $(\Z/p\Z)^*$ is a group), and the map
  $j \mapsto j^{-1}$ is a bijection on $\{1, \ldots, p-1\}$.
  Therefore summing over all $j$ and summing over all $j^{-1}$ gives the same
  sum:
  \[
    H = \sum_{j=1}^{p-1} j^{-1} = \sum_{j=1}^{p-1} j \pmod p.
  \]

  \textbf{Step 2 (Evaluating $\sum_{j=1}^{p-1} j$).}
  \[
    \sum_{j=1}^{p-1} j = \frac{(p-1)p}{2}.
  \]
  Modulo $p$: $p \equiv 0 \pmod p$, so $(p-1)p/2 \equiv 0 \pmod p$.

  \textbf{Conclusion.}
  Therefore $H \equiv \sum_{j=1}^{p-1} j \equiv 0 \pmod p$. \qedhere
\end{proof}

\begin{remark}
  Wolstenholme's stronger theorem (1862) states that for $p \geq 5$:
  \[
    1 + \frac{1}{2} + \cdots + \frac{1}{p-1} \equiv 0 \pmod{p^2}.
  \]
  This is a deeper result and requires analysis modulo $p^2$.
  A prime $p$ is called a \emph{Wilson prime} if $p \mid W_p$, where
  $W_p = \bigl((p-1)!+1\bigr)/p$ is the Wilson quotient.
  As shown in Paper~10 of this series, $W_p \equiv \mu_p + S_p \pmod{p}$,
  where $\mu_p = \bigl((-1)^m(m!)^2+1\bigr)/p$,
  $S_p = \sum_{j=1}^{m} j^{-1} \pmod{p}$, and $m = (p-1)/2$.
  Hence $p$ is a Wilson prime iff $\mu_p + S_p \equiv 0 \pmod{p}$.
  Note that the \emph{full} harmonic sum $\sum_{j=1}^{p-1} j^{-1} \equiv 0 \pmod p$
  always (Theorem~\ref{thm:harmonic}), so only the \emph{half}-sum $S_p$
  carries nontrivial information about Wilson primes.
\end{remark}

% ================================================================
\section{Fermat's Little Theorem via Wilson}
% ================================================================

\begin{theorem}[Fermat's little theorem]\label{thm:fermat}
  Let $p$ be a prime and $a$ an integer with $\gcd(a,p) = 1$.  Then
  \[
    a^{p-1} \;\equiv\; 1 \pmod p.
  \]
\end{theorem}

\begin{proof}[Proof via Wilson]
  Consider the $p-1$ multiples $a, 2a, 3a, \ldots, (p-1)a$.
  We claim these are all distinct modulo $p$ and form a permutation of
  $\{1, 2, \ldots, p-1\}$.

  \textbf{Step 1 (Distinctness).}
  If $ja \equiv ka \pmod p$ with $1 \leq j, k \leq p-1$, then
  $p \mid (j-k)a$.  Since $\gcd(a,p)=1$, we get $p \mid j-k$.
  But $|j-k| < p$, so $j = k$.

  \textbf{Step 2 (They lie in $\{1,\ldots,p-1\}$).}
  Each $ja$ satisfies $\gcd(ja, p) = 1$: since $p$ is prime, $p \nmid j$
  (as $1 \leq j \leq p-1$) and $p \nmid a$ (as $\gcd(a,p)=1$), so $p \nmid ja$,
  giving $ja \not\equiv 0 \pmod p$.

  \textbf{Step 3 (Product comparison).}
  Multiplying all the $p-1$ congruences together:
  \[
    a \cdot 2a \cdot 3a \cdots (p-1)a \;\equiv\; 1 \cdot 2 \cdot 3 \cdots (p-1) \pmod p.
  \]
  The left side is $a^{p-1} \cdot (p-1)!$ and the right side is $(p-1)!$.
  Therefore
  \[
    a^{p-1} \cdot (p-1)! \equiv (p-1)! \pmod p.
  \]

  \textbf{Step 4 (Cancellation using Wilson).}
  By Wilson's theorem, $(p-1)! \equiv -1 \pmod p$, so $\gcd((p-1)!, p) = 1$
  (since $p \nmid (p-1)!+1$ would contradict Wilson; rather Wilson says
  $p \mid (p-1)!+1$, i.e.\ $(p-1)! \equiv -1 \not\equiv 0$).
  Hence $(p-1)!$ is invertible modulo $p$, and we may cancel it:
  \[
    a^{p-1} \equiv 1 \pmod p. \qedhere
  \]
\end{proof}

\begin{remark}
  The standard proof of Fermat's little theorem (using the permutation
  argument in Step 1--3) does not actually require Wilson — the cancellation
  in Step 4 only needs $(p-1)! \not\equiv 0 \pmod p$, which is clear since
  $p$ is prime and none of $1, \ldots, p-1$ is divisible by $p$.
  Wilson's theorem is invoked here to provide the \emph{value} of $(p-1)!$
  and thereby demonstrate the cancellation step more explicitly.
  The proof above shows how Wilson and Fermat are facets of the same
  group-theoretic fact: $(\Z/p\Z)^*$ is a group of order $p-1$.
\end{remark}

\begin{corollary}[Sketch: Euler's theorem via Wilson + structure]\label{cor:euler}
  For $\gcd(a, n) = 1$: $a^{\phi(n)} \equiv 1 \pmod n$.
\end{corollary}

\begin{proof}[Proof sketch]
  By the Chinese Remainder Theorem, $(\Z/n\Z)^* \cong \prod_i (\Z/p_i^{k_i}\Z)^*$
  for the prime factorisation $n = \prod p_i^{k_i}$.
  Each factor $(\Z/p_i^{k_i}\Z)^*$ has order $\phi(p_i^{k_i})$, so
  $a^{\phi(p_i^{k_i})} \equiv 1 \pmod{p_i^{k_i}}$.
  Since $\phi(n) = \prod \phi(p_i^{k_i})$ and $\phi(p_i^{k_i}) \mid \phi(n)$,
  we get $a^{\phi(n)} \equiv 1$ modulo each $p_i^{k_i}$, hence
  $a^{\phi(n)} \equiv 1 \pmod n$ by CRT.
  The role of Wilson-type theorems is to ensure that the unit groups
  $(\Z/p^k\Z)^*$ have order exactly $\phi(p^k) = p^{k-1}(p-1)$, which
  requires the fact that these groups are cyclic (proved using the existence
  of primitive roots, which in turn uses Wilson for $p$).
\end{proof}

% ================================================================
\section{Conclusion}
% ================================================================

Wilson's theorem, while computationally impractical as a primality test,
illuminates deep structural properties of the integers modulo $p$.

The derivation of the first supplement to quadratic reciprocity
($\leg{-1}{p} = (-1)^{(p-1)/2}$) from Wilson via a double counting of $(p-1)!$
is one of the most elegant arguments in elementary number theory.
Wolstenholme's harmonic lemma follows from a simple symmetry argument about
the group $(\Z/p\Z)^*$.  And the proof of Fermat's little theorem via Wilson
exhibits the group-theoretic unity underlying both results.

These four applications demonstrate that Wilson's theorem is not an isolated
curiosity but a foundational tool, connecting primality, quadratic reciprocity,
harmonic sums, and the multiplicative structure of $\Z/p\Z$.

\begin{thebibliography}{9}

\bibitem{Hardy1979}
G.~H.~Hardy and E.~M.~Wright.
\newblock \emph{An Introduction to the Theory of Numbers}, 5th~ed.
\newblock Oxford University Press, 1979.

\bibitem{Ireland1990}
K.~Ireland and M.~Rosen.
\newblock \emph{A Classical Introduction to Modern Number Theory}, 2nd~ed.
\newblock Springer, New York, 1990.

\bibitem{Agrawal2004}
M.~Agrawal, N.~Kayal, and N.~Saxena.
\newblock PRIMES is in P.
\newblock \emph{Ann.\ of Math.}, 160(2):781--793, 2004.

\bibitem{Niven1991}
I.~Niven, H.~S.~Zuckerman, and H.~L.~Montgomery.
\newblock \emph{An Introduction to the Theory of Numbers}, 5th~ed.
\newblock Wiley, New York, 1991.

\bibitem{Crandall1997}
R.~Crandall, K.~Dilcher, and C.~Pomerance.
\newblock A search for Wieferich and Wilson primes.
\newblock \emph{Math.\ Comp.}, 66(217):433--449, 1997.

\end{thebibliography}

\end{document}
